\section{Short paths between representatives}\label{sec:growth}

We retain the optimal path system $\mathcal P$ and the notation
$Q,t,h,L,e,S,F$ from Section~\ref{sec:paths}, and assume throughout
this section that $t\ge 9$.
We shall choose one vertex of $S$ for each vertex of $Q$ so that an
outside path between two chosen vertices can be inserted into the
system without losing any of its old interior vertices.  Neighborhood
expansion will then force such an outside path with at most six
internal vertices.

Orient each chain of the forest represented by $\mathcal P$.
For each \emph{target} $q\in Q$, define its \emph{representative}
$\rho(q)\in S$ as follows.  If $q$ is the first vertex of its chain,
including a singleton chain, set $\rho(q)=q$.  Otherwise its incoming
path, read in the chosen orientation, has the form
\[
 p,z_1,z_2,\ldots,z_a,q,\qquad a\ge 1;
\]
set $\rho(q)=z_1$, the first interior vertex after the predecessor $p$.
The segment from $\rho(q)$ to $q$ therefore contains all $a$ interior
vertices of the incoming path.  The representatives are distinct,
since different incoming paths have disjoint interiors and only
chain starts are represented by clique vertices.  Write
$R=\{\rho(q):q\in Q\}$, so $|R|=t$.

\begin{lemma}[Replacing incoming paths]\label{grow:replacement}
Fix an orientation of the chains and distinct targets $q,q'\in Q$.
Suppose an outside path relative to $S$ joins $\rho(q)$ to $\rho(q')$
and has $d\ge 1$ internal vertices.  Let $r\in\{0,1,2\}$ be the number
of these two targets that have an incoming path.
Deleting those $r$ paths from $\mathcal P$ and inserting one new path
from $q$ to $q'$ gives a valid path system with total amount $L+d$
and $e+1-r$ paths.  It retains all the interior vertices of
$\mathcal P$.
\end{lemma}

\begin{proof}
At each target with an incoming path, retain its segment from the
representative to the target.  At a chain start the corresponding
segment consists only of the target.  Concatenate the reverse of the
segment at $q$, the given outside path, and the segment at $q'$.
Figure~\ref{fig:representatives} illustrates the construction with two
incoming paths.
The two retained segments have disjoint interiors, and the outside
path has its interior in $F=G-S$, so the resulting path is simple.
In particular, a retained segment contains no clique vertex except
its own target.  This remains true when the two targets are
consecutive on one chain.

After deleting the incoming paths, each target has degree at most
one in the represented forest, and the targets lie in different
components.  If they belonged to different chains, this is immediate.
If they belonged to the same chain, deleting the incoming edge at
the later target separates them.  The new edge between the targets
thus creates neither a cycle nor a vertex of degree greater than
two.  It gives a valid linear forest.

If the deleted paths have total amount $A$, their retained segments
contribute exactly $A$ interior vertices to the new path.  Its amount
is therefore $A+d$, while the unchanged paths contribute $L-A$.
The total is $L+d$, and the number of paths is $e-r+1$.
\end{proof}

\begin{lemma}[One-vertex detours]\label{grow:one}
For every orientation of the chains, no outside path relative to $S$
with one internal vertex joins two distinct representatives.
Nor can such a path join consecutive vertices of any path of
$\mathcal P$.
\end{lemma}

\begin{proof}
The first type of path would give a system of amount $L+1$ by
Lemma~\ref{grow:replacement}.  For the second type, replace the step
between the consecutive vertices by the proposed two-edge path.
This keeps the represented forest and every old interior vertex,
again increasing the amount to $L+1$.

Since $L<h$, either $L+1<h$, contrary to maximality of $L$, or
$L+1=h$.  In the latter case the new system has at most $t$ incident
clique vertices, and hence lies in the forbidden interval of
Lemma~\ref{path:interval}.
\end{proof}

We next translate these exclusions into growth of the balls defined
in Section~\ref{sec:paths}.  Fix an orientation and a representative
$x\in R$.  An outside neighbor $u\in B_0(x;S)$ exists because
$|S|\le k$ and $\delta(G)\ge k$.  By Lemma~\ref{grow:one}, $u$ misses
the other $t-1$ representatives.  Since its closed neighborhood in
$F$ lies in $B_1(x;S)$, we obtain
\begin{equation}\label{grow:degree}
 |B_1(x;S)|
 \ge 1+\deg_F(u)
 \ge 1+k-\bigl(|S|-(t-1)\bigr)
 =k-|S|+t.
\end{equation}
Moreover, if some $y\in S\setminus R$ is missed by every vertex of
$B_0(x;S)$, the same count improves the right-hand side by one.
Only the one-vertex exclusions are needed for this degree estimate.

\begin{lemma}[Independence in the second ball]\label{grow:weights}
Fix an orientation and $x\in R$.  Suppose no outside path relative
to $S$ with one or two internal vertices joins $x$ to any other
representative.  Then
\[
 \alpha\bigl(B_2(x;S)\bigr)\ge 2.
\]
If, in addition, $B_1(x;S)$ is not a clique and $t<k$, then
$\alpha(B_2(x;S))\ge 3$.
\end{lemma}

\begin{proof}
Every vertex of $B_1(x;S)$ misses $R\setminus\{x\}$: an adjacency
to another representative would give an outside path with one or
two internal vertices.

Suppose first that $B_2(x;S)$ were a clique.  It contains
$B_1(x;S)$, whose order is at least $k-|S|+t\ge t$ by
\eqref{grow:degree}.  Since the maximum clique size is $t$, this forces
\[
 B_1(x;S)=B_2(x;S),\qquad |B_1(x;S)|=t,\qquad |S|=k.
\]
Each vertex $z$ of this ball has at most $t-1$ neighbors in $F$,
because all those neighbors lie in $B_2(x;S)$.  It also has at most
$k-t+1$ neighbors in $S$, since it misses the other representatives.
As $\deg_G(z)\ge k$, both bounds are attained.  In particular every
such $z$ is adjacent to $x$.  The ball together with $x$ is then a
clique of order $t+1$, a contradiction.  Thus $B_2(x;S)$ is not a
clique and has an independent pair.

Now suppose that $B_1(x;S)$ contains an independent pair $I$.
Its closed neighborhood in $F$ lies in $B_2(x;S)$, and its
neighborhood in $S$ avoids $R\setminus\{x\}$.  The expansion bound
of Lemma~\ref{red:expansion} gives
\begin{equation}\label{grow:pair-expansion}
 |B_2(x;S)|
 \ge 2k+1-\bigl(|S|-(t-1)\bigr)
 =2k-|S|+t
 \ge k+t.
\end{equation}
When $t<k$, this is greater than $\max(k,2t)$, so
Lemma~\ref{size:test} gives an independent triple.
\end{proof}

The independent pairs in the second balls may already give a
contradiction.  When they do not, the only remaining obstruction to
obtaining independent triples is a first ball that is a clique of
order $t$ with $|S|=k$.  In that case, the structure of the path
system will supply an additional excluded neighbor for at least two
representatives.

\begin{lemma}[A short outside path]\label{grow:short}
Suppose $t\ge 9$.  For every fixed orientation of the chains of
$\mathcal P$, two distinct representatives are joined by an outside
path relative to $S$ with $1\le d\le 6$ internal vertices.
By Lemma~\ref{grow:one}, such a path in fact has $2\le d\le 6$.
\end{lemma}

\begin{proof}
Fix an orientation and suppose that no such path exists.
Lemma~\ref{path:separation} makes the $t$ balls
$B_2(x;S)$, $x\in R$, pairwise disjoint and anticomplete.
Each has an independent pair by Lemma~\ref{grow:weights}, so together
they give an independent set of order at least $2t$.
If $2t>k$, this contradicts $\alpha(G)\le k$; this also settles the
case $t=k$ without using the independent-triple conclusion.

We may therefore assume $2t\le k$.  The bound
$\lfloor k/2\rfloor\le t$ from Theorem~\ref{clq:bound} now gives
\begin{equation}\label{grow:parity}
 k\in\{2t,2t+1\}.
\end{equation}
In particular $t<k$, and it suffices to find two representatives
whose first balls are not cliques.  Lemma~\ref{grow:weights} would
then give a combined independent set of order at least
$2t+2>k$.

If $|S|<k$, equation~\eqref{grow:degree} gives
$|B_1(x;S)|>t$ at every representative, so any two suffice.
It remains to consider $|S|=k$.  In this case
$L=k-t>0$, and hence $e\ge 1$.

Suppose first that $e\ge 2$.  At least two representatives are
interior vertices of paths of $\mathcal P$.  For each such
representative $x$, let $y$ be the next vertex toward its target.
This vertex is not in $R$: it is either a later interior vertex of
the same path, or, when the amount is one, the target itself, which
is not a chain start.  By the consecutive-vertex assertion of
Lemma~\ref{grow:one}, every outside neighbor of $x$ misses $y$.
The improved degree count following \eqref{grow:degree} yields
\[
 |B_1(x;S)|\ge k-|S|+t+1=t+1.
\]
Thus the two chosen first balls are not cliques.

Finally suppose that $e=1$.  The represented forest has one
nonsingleton chain, consisting of a single edge, and $t-2\ge 7$
singleton chains.  Let $x$ be the representative of any singleton
chain.  Reverse the nonsingleton chain, leaving every other
orientation fixed, and let $y$ be its new starting clique vertex.
The vertex $y$ was not in the original representative set $R$,
whereas $x$ is a representative in both orientations.
Lemma~\ref{grow:one}, applied to the reversed orientation, shows that
every outside neighbor of $x$ misses $y$.  Consequently
$|B_1(x;S)|\ge t+1$ once again.  Choosing two singleton
representatives gives the two required noncliques and the final
contradiction.  The reversed orientation was used only for the
unconditional one-vertex exclusion; the assumed absence of paths
with up to six internal vertices concerned the original orientation
alone.
\end{proof}
